\section{Cardinal interpolation and Fourier transformation}
\label{sec:fourier}

The auxiliary pair will be assembled from two summable lists of interpolation
data.  We first realize each list by an entire function whose values and first
derivatives are determined by its coordinates.  Gaussian damping then turns this function
into a radial Schwartz function, and Fourier transformation gives a second
kernel with rapidly decreasing jets at distant nodes.  The estimates in this
section apply to arbitrary summable lists, before any finite approximation is
introduced.

The spectral construction adapts the cardinal-measure and Gaussian-transform
method of OpenAI~\citep[Section~2]{SharpPlanar2026}.  The formulas and estimates needed
for the present target are proved below; no packing sign bound is used.

\subsection{Shells and interpolation nodes}

Set the construction parameters and radial coordinate as follows:
\begin{equation}\label{eq:normalization}
 b=\frac{\sqrt3}{2},\qquad h=\frac25,\qquad
 B=\frac43=\frac1{b^2},\qquad s=b|x|^2.
\end{equation}
The same radial coordinate is used on the Fourier side.  For the unit-covolume
triangular lattice \(A\), take the basis
\[
 u=\sqrt{2/\sqrt3}\,(1,0),\qquad
 v=\sqrt{2/\sqrt3}\,(1/2,\sqrt3/2).
\]
Then
\[
 b|ju+\ell v|^2=j^2+j\ell+\ell^2\qquad(j,\ell\in\Z).
\]
Modulo~3 the right side equals \((j-\ell)^2\), and is therefore 0 or 1.
Modulo~2 it is even only when both \(j\) and \(\ell\) are even, in which
case it is divisible by~4.  These two restrictions give the residue classes
listed below.  Thus every nonzero shell belongs to the positive part of the
following periodic set:
\begin{equation}\label{eq:node-set}
 \begin{gathered}
 L=\{0,1,3,4,7,9\},\qquad \mathcal Z=12\Z+L,\\
 \mathcal N=\mathcal Z\cap(0,\infty),\qquad
 \mathcal N_0=\{0\}\cup\mathcal N.
 \end{gathered}
\end{equation}
The inclusion is strict.  For example, \(24\in\mathcal N\), but if
\(j^2+j\ell+\ell^2=24\), the parity implication makes \(j,\ell\) even.
Dividing by four gives the same quadratic form with value 6, so its two
arguments are again even.  The original value would then be divisible by 16,
contrary to \(16\nmid24\).  Such additional nodes will also carry prescribed
interpolation data.

Let \(J\) denote rotation through a positive right angle.  Since
\(\det(u,v)=1\), the vectors \(-Jv,Ju\) are dual to \(u,v\).  They generate
\(JA\), because \(-v,u\) form another basis of \(A\).  Hence \(A^*=JA\), so
the direct and dual lattices have the same radial shells.

\subsection{The normalized Gaussian theorem}

The Fourier transform uses the convention fixed in the introduction.
The common direct and dual shell set allows us to impose contacts at every
positive node of $\mathcal N$.  The following theorem is the auxiliary
result proved in Sections~\ref{sec:fourier}--\ref{sec:signs}.

\begin{theorem}[Normalized Gaussian Fourier pair]\label{thm:normalized-pair}
For every finite real $k\ge2.36$, put $T_k(s)=k^{-1}e^{-k(s-1)}$.
There are entire functions $H_1,H_2$, real on the real axis, and a real radial
Schwartz function $f$ on $\R^2$ such that
\[
 f(x)=e^{-\pi h b\abs x^2}H_1(b\abs x^2),\qquad
 \widehat f(\xi)=e^{-\pi h b\abs\xi^2}H_2(b\abs\xi^2).
\]
For every real $s\ge0$,
\begin{equation}\label{eq:normalized-signs}
 H_1(s)\le T_k(s),\qquad H_2(s)\ge0.
\end{equation}
At every $n\in\mathcal N$, the values and first derivatives satisfy
\begin{equation}\label{eq:normalized-jets}
 \begin{gathered}
 H_1(n)=T_k(n),\qquad H_1'(n)=T_k'(n)=-e^{-k(n-1)},\\
 H_2(n)=H_2'(n)=0.
 \end{gathered}
\end{equation}
\end{theorem}

The full parameter range $k\ge2.36$ and every artificial node are part of
the assertion, beyond the contacts needed for the energy inequalities.
The theorem asserts existence; uniqueness will be proved only in the fixed
coefficient system of Proposition~\ref{prop:exact-interpolation}.

For $\alpha\ge1$, multiplying the pair by $ke^{-k}$ with
$k=\pi(\alpha/b-h)$ gives a minorant of $e^{-\pi\alpha\abs x^2}$.
A Fourier-complement construction supplies $0<\alpha<1$.
Section~\ref{sec:energy} proves this conversion and the sharp Gaussian energy
comparison; Section~\ref{sec:mixtures} then proves
Theorem~\ref{thm:universal}.

The atomic companion \citep[Sections~3--4]{AtomicTriangular2026} uses finite
atomic columns, a prescribed node set omitting $24$, and Gaussian damping
$h=17/50$.  Here $h=2/5$, and both jets are imposed at every positive node
of the periodic superset.  The common energy conclusion therefore does not
identify the two auxiliary-function theorems.

\subsection{The sine product and its node coordinates}

We now begin the construction of the theorem's pair.  The first ingredient
is a function with a double zero at every prescribed node.  Define
\begin{equation}\label{eq:periodic-polynomial}
 P(s)=\prod_{a\in L}\left(2\sin\frac{\pi(s-a)}{12}\right)^2
     =\sum_{l=-6}^{6}P_l e^{i\pi t_l s},\qquad t_l=\frac l6.
\end{equation}
This is an entire trigonometric polynomial, nonnegative and 12-periodic on the
real axis, with a double zero at each point of \(\mathcal Z\) and no other real
zeros.  Its Fourier coefficients are
\begin{equation}\label{eq:periodic-coefficients}
 (P_0,P_1,\ldots,P_6)
 =\left(5,-1+2ib,1+2ib,-2,-\frac12+ib,-1-2ib,1\right),
 \qquad P_{-l}=\overline{P_l}.
\end{equation}
Indeed, setting \(w=e^{i\pi s/6}\) expresses the product as
\(w^{-6}\prod_{a\in L}(w-e^{i\pi a/6})^2\).  The scalar factor is one because
\(\sum_{a\in L}a=24\); multiplying this degree-twelve polynomial gives
\eqref{eq:periodic-coefficients}.

At a node \(n\in\mathcal N_0\), write
\begin{equation}\label{eq:node-expansion}
 P(n+u)=Q_nu^2\bigl(1+D_nu+O(u^2)\bigr).
\end{equation}
The periodicity of \(P\) makes \(Q_n,D_n\) depend only on the residue \(a\) of \(n\):
\begin{equation}\label{eq:node-constants}
\begin{array}{c|cc}
 a&Q_n&D_n\\ \hline
 0&\pi^2/3&-7\pi\sqrt3/18\\
 1&(2\pi^2/3)(2-\sqrt3)&\pi(3+\sqrt3)/18\\
 3&(2\pi^2/3)(2-\sqrt3)&-\pi(3+\sqrt3)/18\\
 4&\pi^2/3&7\pi\sqrt3/18\\
 7&(2\pi^2/3)(2+\sqrt3)&-\pi(3-\sqrt3)/18\\
 9&(2\pi^2/3)(2+\sqrt3)&\pi(3-\sqrt3)/18
\end{array}
\end{equation}
To obtain the table, remove the vanishing sine factor and take the value and
logarithmic derivative of the remaining product:
\[
 Q_a=\left(\frac\pi6\right)^2
       \prod_{d\in L\setminus\{a\}}
        \left(2\sin\frac{\pi(a-d)}{12}\right)^2,
 \qquad
 D_a=\frac\pi6\sum_{d\in L\setminus\{a\}}
             \cot\frac{\pi(a-d)}{12}.
\]
The angle-addition formulas give the displayed entries.  In particular,
\begin{equation}\label{eq:node-constant-bounds}
 Q_n>\frac74=1.75,\qquad |D_n|<\frac{53}{25}=2.12.
\end{equation}
For these strict comparisons, \(157/50<\pi<22/7\) and
\(\sqrt3<1733/1000\) suffice.  The smallest \(Q_n\) is greater than
\((2/3)(157/50)^2(267/1000)>7/4\), while the largest \(|D_n|\) is
\(7\pi\sqrt3/18<(22/18)(1733/1000)<53/25\).
The upper bound for \(\pi\) is proved in Section~\ref{sec:finite}; the lower
bound for \(\pi\) and the upper bound for \(\sqrt3\) are proved in
Appendix~\ref{sec:arithmetic}.

\subsection{A spectral realization of summable cardinal data}

Double-pole value terms and simple-pole derivative terms are familiar from
the one-dimensional constructions of
Cohn and Kumar~\citep[Proposition~9.6]{CohnKumar2007} and Carneiro et~al.~\citep[Section~1, (1.2)--(1.3)]{CarneiroLittmannVaaler2013}.
Here we need their planar Fourier transform and bounds uniform in the node
location.  A compactly supported spectral measure provides both.

A coefficient list consists of a real constant \(C\) and real pairs
\((c_n,d_n)\) for \(n\in\mathcal N_0\), with the unweighted norm
\[
 \|a\|_1:=|C|+\sum_{n\in\mathcal N_0}(|c_n|+|d_n|)<\infty.
\]
The letter \(a\) denotes the full list.  Away from the nodes, put
\begin{equation}\label{eq:cardinal-function}
 p_a(s)=P(s)\left(C+
       \sum_{n\in\mathcal N_0}
          \left(\frac{c_n}{(s-n)^2}+\frac{d_n}{s-n}\right)\right).
\end{equation}
We suppress the subscript when the list is understood.

\begin{lemma}[Cardinal measure]\label{lem:cardinal-measure}
The function in \eqref{eq:cardinal-function} extends to an entire function,
real on the real axis, with jets
\begin{equation}\label{eq:cardinal-jets}
 p_a(n)=Q_nc_n,\qquad p_a'(n)=Q_n(d_n+D_nc_n)
       \quad(n\in\mathcal N_0).
\end{equation}
There is a finite complex measure \(\mu_a\) on \([-1,1]\) such that
\begin{equation}\label{eq:full-measure}
 p_a(s)=\int_{-1}^{1}e^{i\pi ts}\,\dd\mu_a(t)\qquad(s\in\C),
\end{equation}
\(\mu_a(-E)=\overline{\mu_a(E)}\) for every Borel set \(E\), and
\begin{equation}\label{eq:spectral-variation}
 \|\mu_a\|_{\mathrm{TV}}
 \le 25|C|+36\sum_{n\in\mathcal N_0}(|c_n|+|d_n|).
\end{equation}
For real \(s\), the same value can be evaluated from a folded measure on
\([0,1]\):
\begin{equation}\label{eq:folded-evaluation}
 p_a(s)=\Rea\int_0^1e^{i\pi ts}\,\dd\nu_a(t).
\end{equation}
The atoms of \(\nu_a\) are \(Ce_jP_j\) at \(t_j\), where \(e_0=1\) and
\(e_j=2\) for \(1\le j\le6\).  Its density on the open frequency intervals is
\begin{equation}\label{eq:folded-density}
 d_j(t)=2\pi\sum_{n\in\mathcal N_0}e^{-i\pi tn}
       \sum_{l=j}^{6}
           \bigl[-\pi c_n(t_l-t)+id_n\bigr]P_le^{i\pi t_ln}
 \quad(t_{j-1}<t<t_j).
\end{equation}
The same variation bound holds for \(\nu_a\).
\end{lemma}

\begin{proof}
First consider one node \(n\).  For \(u=s-n\) and real \(t_l\), direct
integration gives the entire identities
\begin{align*}
 \frac{e^{i\pi t_lu}-1}{u}
    &=i\pi\int_0^{t_l}e^{i\pi tu}\,\dd t,\\
 \frac{e^{i\pi t_lu}-1-i\pi t_lu}{u^2}
    &=-\pi^2\int_0^{t_l}(t_l-t)e^{i\pi tu}\,\dd t.
\end{align*}
The integrals include the removable values at \(u=0\).  Multiply these
identities by \(P_le^{i\pi t_ln}\) and sum over \(l\).  The subtracted
constant and linear terms cancel because \(P(n)=P'(n)=0\).  We therefore
obtain spectral measures for the two entire node fractions in
\eqref{eq:cardinal-function}.  Their negative-frequency parts are conjugate
reflections of their positive-frequency parts.  Folding doubles the latter
parts and gives exactly \eqref{eq:folded-density}.  These node measures are
absolutely continuous; the atoms come only from \(CP\), with the zero
frequency counted once and every positive frequency counted twice.

The coefficient row in \eqref{eq:periodic-coefficients} gives
\((|P_1|,\ldots,|P_6|)=(2,2,2,1,2,1)\).  The constant column consequently has
variation \(|P_0|+2\sum_{l=1}^6|P_l|=25\).  Taking absolute values in the two
integrals above gives, per unit \(c_n\) and \(d_n\), respectively,
\[
 \pi^2\sum_{l=1}^6t_l^2|P_l|=\frac{65\pi^2}{18}<36,
 \qquad
 2\pi\sum_{l=1}^6t_l|P_l|=\frac{32\pi}{3}<36.
\]
The phases involving \(n\) have modulus one at real frequencies, so these
bounds are independent of the node.  They apply to the full measure and to its
folded form.

For an absolutely summable list, the node measures now converge in total
variation, proving \eqref{eq:spectral-variation}.  On every compact subset
of the complex \(s\)-plane, \(e^{i\pi ts}\) and each of its \(s\)-derivatives
are uniformly bounded for \(|t|\le1\).  Thus the measure integral and all its
derivatives are locally uniform limits of those for finite lists.  On a
compact set avoiding the nodes, the reciprocal factors in the cardinal sum
are uniformly bounded over all \(n\), so absolute summability also gives
uniform convergence of that sum there.  The integral thus defines its entire
extension, with the stated removals.

For a finite list, \eqref{eq:node-expansion} gives
\eqref{eq:cardinal-jets}: the designated double and simple fractions give the
two displayed terms, while every other term vanishes to order at least two at
the node.  Locally uniform convergence of the functions and their derivatives
passes these identities to the full list.  Finally, conjugate symmetry of the
measure gives reality on the real axis and the folded formula.
\end{proof}

The entire extension is always the full complex integral
\eqref{eq:full-measure}.  The real part in \eqref{eq:folded-evaluation} is an
evaluation rule for real \(s\), not a definition at complex arguments.

\subsection{Gaussian damping and the transformed kernel}

For \(-1\le t\le1\), set
\begin{equation}\label{eq:transformed-kernel}
 z(t)=-\frac{B}{t+ih}-ih,\qquad
 \lambda(t)=\frac{i}{b(t+ih)}.
\end{equation}
For a coefficient list \(a\), define
\begin{equation}\label{eq:transformed-function}
 K_{p_a}(s)=\int_{-1}^1\lambda(t)e^{i\pi z(t)s}\,\dd\mu_a(t).
\end{equation}
The image of \([-1,1]\) under \(z\) is compact and \(\lambda\) is bounded
there, so the same compact-uniform argument makes \(K_{p_a}\) entire.  Also,
\[
 \lambda(-t)=\overline{\lambda(t)},\qquad
 z(-t)=-\overline{z(t)}.
\]
It follows that \(K_{p_a}\) is real on the real axis, where
\begin{equation}\label{eq:folded-transformed-function}
 K_{p_a}(s)=\Rea\int_0^1
                   \lambda(t)e^{i\pi z(t)s}\,\dd\nu_a(t).
\end{equation}

\begin{remark}[Pointwise frequency consistency]
The rational formulas in \eqref{eq:transformed-kernel} also make sense for
complex \(t\ne-ih\).  Since \(z(t)+ih=-B/(t+ih)\), the image avoids \(-ih\),
and substitution twice gives
\[
 z(z(t))=t,\qquad
 \lambda(t)\lambda(z(t))
 =\frac{i}{b(t+ih)}\frac{-i(t+ih)}{bB}=1.
\]
These are pointwise frequency and amplitude identities for the Gaussian
transform.  They do not define a composition of the folded real-contour
integral: \(z\) does not preserve that contour, and a transformed profile has
not thereby been given a cardinal measure on \([-1,1]\).  The involution for
the radial functions follows from the genuine Fourier transform below.
\end{remark}

\begin{lemma}[Fourier pairs]\label{lem:fourier-pair}
For every coefficient list \(a\), the radial functions
\(e^{-\pi hs}p_a(s)\) and \(e^{-\pi hs}K_{p_a}(s)\) are real Schwartz
functions on \(\R^2\) and are Fourier transforms of one another.
Consequently, for any two lists, the entire functions
\begin{equation}\label{eq:coupled-pair}
 H_1=p_1+K_{p_2},\qquad H_2=p_2+K_{p_1}
\end{equation}
give the real radial Schwartz Fourier pair
\[
 f(x)=e^{-\pi hb|x|^2}H_1(b|x|^2),\qquad
 \widehat f(\xi)=e^{-\pi hb|\xi|^2}H_2(b|\xi|^2).
\]
\end{lemma}

\begin{proof}
For a complex parameter \(\gamma\) with \(\Rea\gamma>0\), the Gaussian
transform in dimension two, with kernel \(e^{-2\pi i x\cdot\xi}\), is
\begin{equation}\label{eq:complex-gaussian-transform}
 \widehat{e^{-\pi\gamma|x|^2}}(\xi)
       =\gamma^{-1}e^{-\pi|\xi|^2/\gamma}.
\end{equation}
For real \(\gamma>0\), squaring the Gaussian integral and using polar
coordinates gives the one-dimensional value \(\gamma^{-1/2}\) at frequency
zero.  Integration by parts shows that the one-dimensional transform obeys
\(F'(\xi)=-(2\pi\xi/\gamma)F(\xi)\).  Solving this differential equation and
taking the product in two coordinates proves
\eqref{eq:complex-gaussian-transform} for positive real \(\gamma\).  Both
sides are holomorphic on the right half-plane; differentiation under the
integral is locally justified by a polynomial times a decaying Gaussian.
The identity theorem gives the complex case.

After multiplication by \(e^{-\pi hs}\), the integrand in
\eqref{eq:full-measure} is a Gaussian with parameter
\(\gamma=b(h-it)=-ib(t+ih)\), whose real part is the fixed positive number
\(bh\).  Using \(B=1/b^2\), its transform is
\[
 \frac1{b(h-it)}
    \exp\left(-\frac{\pi|\xi|^2}{b(h-it)}\right)
 =e^{-\pi hs}\lambda(t)e^{i\pi z(t)s},\qquad s=b|\xi|^2.
\]
Integrating this identity against \(\mu_a\) proves the forward Fourier
formula.  The exchanges with Fourier integration, derivatives, and polynomial
weights are justified uniformly: the measure is finite, the frequency
interval is compact, and every differentiated real-space Gaussian is bounded
by a fixed polynomial times a decaying Gaussian.  On the transformed side,
the real part of the Gaussian parameter is \(h/[b(h^2+t^2)]\), with a positive
lower bound on \([-1,1]\).  The same domination applies there.  Both functions
are therefore Schwartz.

Fourier transformation twice reflects a Schwartz function.  The functions
here are radial, so reflection fixes them and proves the reverse transform.
Linearity gives \eqref{eq:coupled-pair}.
\end{proof}

\subsection{The jet operator and its distant outputs}

For a function differentiable at the positive nodes, define the jet projection
\begin{equation}\label{eq:jet-projection}
 (\mathcal Ju)_{c,n}=\frac{u(n)}{Q_n},\qquad
 (\mathcal Ju)_{d,n}=\frac{u'(n)-D_nu(n)}{Q_n}
       \quad(n\in\mathcal N).
\end{equation}
By Lemma~\ref{lem:cardinal-measure}, \(\mathcal Jp_a\) is precisely the
positive-node part of \(a\).  The three remaining coordinates, in the order
\((c_0,d_0,C)\), are called the extra coordinates.  They make no direct
contribution to the jets at positive nodes; we will choose them in
Section~\ref{sec:finite}.  Put
\begin{equation}\label{eq:jet-operator}
 Sa=-\mathcal J K_{p_a}.
\end{equation}
For two full lists \(a_1,a_2\), let \(a_i^+\) denote their positive-node
parts.  The coupled pair \eqref{eq:coupled-pair} then satisfies
\[
 \mathcal JH_1=a_1^+-Sa_2,\qquad
 \mathcal JH_2=a_2^+-Sa_1.
\]
Thus the interpolation problem asks for these two expressions to equal
\(\mathcal JT_k\) and zero, respectively.  This explains the minus sign
in the definition of \(S\).  The next estimates bound the transformed
contribution, especially at distant output nodes, uniformly in the locations
of the input coefficients.

The same letter \(S\) denotes its restriction to positive-node inputs with
zero extras.  Subscripts specify input and output restrictions.  All list
norms are unweighted sums of absolute values, including the sum of the norms
on a pair of lists; operator norms are induced by them.  A finite matrix
therefore has the maximum absolute column-sum norm.

\begin{lemma}[Kernel and tail bounds]\label{lem:kernel-bounds}
For \(0\le t\le1\),
\begin{equation}\label{eq:global-kernel-bounds}
 |\lambda(t)|<3,\qquad |z(t)|\le\frac{44}{15}<3,
 \qquad \pi\Ima z(t)>.187.
\end{equation}
The last quantity exceeds \(.70\) when \(0\le t\le5/6\).  On
\(5/6\le t\le1\), the sharper bounds are
\begin{equation}\label{eq:last-interval-bounds}
 |\lambda(t)|<1.25,\qquad |z(t)|<1.32,
 \qquad \pi\Ima z(t)>.187+2.48(1-t).
\end{equation}
The operator \(S\), with extra input coordinates included if desired, maps
coefficient lists boundedly to positive-node \(\ell^1\) lists and has norm
less than 4000.  More precisely, if
\[
 w\ge\sum_{n\in\mathcal N_0}(|c_n|+|d_n|),\qquad c\ge|C|,
\]
then for every positive integer \(r\),
\begin{equation}\label{eq:tail-bound}
 \sum_{\substack{n\in\mathcal N\\n\ge r}}
       \bigl(|(Sa)_{c,n}|+|(Sa)_{d,n}|\bigr)
 \le U_r(w,c):=
 800(w+c)\frac{e^{-.70r}}{1-e^{-.70}}
 +5.3\left(\frac{2.54w}{r}+2c\right)
                  \frac{e^{-.187r}}{1-e^{-.187}}.
\end{equation}
For several lists, the same formula bounds the sum of their tail outputs when
\(w,c\) bound the corresponding sums of input coefficients.
\end{lemma}

\begin{proof}
Write \(q(t)=\pi\Ima z(t)\) within this proof.  Direct calculation gives
\begin{align*}
 |\lambda(t)|&=\frac1{b\sqrt{t^2+h^2}},\\
 |z(t)|^2&=h^2+\frac{B^2-2Bh^2}{t^2+h^2},\\
 q(t)&=\pi h\left(\frac{B}{t^2+h^2}-1\right).
\end{align*}
Since \(B>2h^2\), all three expressions decrease for \(t\ge0\).  At zero,
\(|\lambda|=5/\sqrt3<3\) and \(|z|=44/15\).  At the other relevant endpoints,
\[
 q(1)=\frac{26\pi}{435}>.187,\qquad
 q(5/6)=\frac{862\pi}{3845}>.70,
\]
and
\[
 |\lambda(5/6)|=\frac{60}{\sqrt{2307}}<\frac54,
 \qquad
 |z(5/6)|^2=\frac{33476}{19225}<\left(\frac{33}{25}\right)^2.
\]
Finally,
\[
 -q'(t)=\frac{2\pi hBt}{(t^2+h^2)^2}.
\]
The function \(t/(t^2+h^2)^2\) decreases for \(t>h/\sqrt3\).  Its value at
the last endpoint therefore gives \(-q'(t)\ge2000\pi/2523>2.48\) throughout
\([5/6,1]\).  Integrating from \(t\) to 1 proves the last inequality in
\eqref{eq:last-interval-bounds}.  All strict numerical comparisons here
follow from \(157/50<\pi<22/7\) and squaring the positive rational bounds.

One unit of spectral variation at frequency \(t\) contributes at most
\begin{equation}\label{eq:jet-kernel-cost}
 \frac{|\lambda(t)|}{Q_n}
       \bigl(1+|D_n|+\pi|z(t)|\bigr)e^{-q(t)n}
\end{equation}
to the sum of the two absolute jet coordinates at \(n\).  By
\eqref{eq:node-constant-bounds}, its prefactor is less than 22 globally and
less than 5.3 on the last interval, since
\[
 \frac{3(1+2.12+3\pi)}{1.75}<22,
 \qquad
 \frac{1.25(1+2.12+1.32\pi)}{1.75}<5.3.
\]
The variation bound is at most \(36\|a\|_1\).  Summing
\eqref{eq:jet-kernel-cost} over all positive integers, a superset of the
positive nodes, therefore gives
\[
 \|S\|\le\frac{22\cdot36}{e^{.187}-1}<4000.
\]
For the final comparison, \(e^{.187}-1>.187+(.187)^2/2>.198\).

For the sharper tail, split the folded measure into \([0,5/6]\) and
\((5/6,1]\), assigning the atom at \(5/6\) to the first part.  The first part
has variation at most \(36w+25c\).  Its contribution to the row pair at \(n\)
is therefore at most
\[
 22(36w+25c)e^{-.70n}\le800(w+c)e^{-.70n}.
\]
On the last open frequency interval, only \(l=6\) remains in
\eqref{eq:folded-density}.  Here \(|P_6|=1\) and
\(\pi(1-t)\le\pi/6<1\), so the density has absolute value at most \(2\pi w\).
The only atom in \((5/6,1]\) is the one at 1, of mass \(2C\).  This part
contributes at most
\[
 5.3e^{-.187n}
       \left(2\pi w\int_{5/6}^1e^{-2.48n(1-t)}\,\dd t+2c\right)
 \le5.3e^{-.187n}\left(\frac{2.54w}{n}+2c\right),
\]
because \(2\pi/2.48<2.54\).  Summing over the integers \(n\ge r\) and using
\(1/n\le1/r\) proves \eqref{eq:tail-bound}.  Only absolute-value sums were
used, so the same argument applies to several lists with their summed inputs.
\end{proof}

The same estimates control every derivative.  If
\(V=25|C|+36\sum_n(|c_n|+|d_n|)\), then for every integer \(j\ge0\) and real
\(s\ge0\), differentiation of the spectral integrals gives
\begin{equation}\label{eq:uniform-derivative-bounds}
 |p_a^{(j)}(s)|\le\pi^jV,\qquad
 |K_{p_a}^{(j)}(s)|\le3(3\pi)^j e^{-.187s}V.
\end{equation}
Keeping the two frequency ranges separate as above gives, for \(s>0\),
\begin{equation}\label{eq:refined-derivative-bound}
 |K_{p_a}^{(j)}(s)|
 \le3(3\pi)^j(36w+25c)e^{-.70s}
    +1.25(1.32\pi)^j
          \left(\frac{2.54w}{s}+2c\right)e^{-.187s}.
\end{equation}
In particular, the jet, tail, and derivative bounds are uniform in the input
node locations and require no finite-support assumption.
